\chapter{Data Factory Pipelines}
\index{Data Factory}\index{pipelines}\index{Copy activity}\index{Lookup activity}\index{ForEach}\index{control flow}\index{parameters}\index{dynamic expressions}\index{data contracts}\index{metadata-driven ingestion}%

\begin{objectives}
\begin{itemize}
    \item Build Fabric pipelines using the Copy activity, Lookup, and Get Metadata.
    \item Compose control flow with ForEach, If Condition, Until, and Switch.
    \item Parameterize pipelines with parameters, variables, and dynamic expressions.
    \item Enforce data contracts: validate expected schema before writing to Silver, and classify breaking versus non-breaking changes.
    \item Build a metadata-driven pipeline that copies a list of tables from Azure SQL Database to OneLake with a Lookup and a ForEach loop.
    \item Design an ingestion framework that onboards hundreds of tables without manual pipeline creation.
\end{itemize}
\end{objectives}

\begin{crosscloud}
The visual, activity-based orchestrator with a rich connector library is a pattern every cloud offers. Fabric Pipelines descends directly from Azure Data Factory, so the mapping is especially tight in the Microsoft column.

\vspace{2mm}
{\footnotesize
\begin{tabular}{@{}p{2.8cm}p{3.2cm}p{3.2cm}p{3.2cm}p{3.2cm}@{}}
\toprule
\textbf{Capability} & \textbf{Fabric} & \textbf{AWS} & \textbf{GCP} & \textbf{Snowflake} \\
\midrule
Visual orchestrator & Data Factory pipelines & Glue Workflows / Step Functions & Cloud Composer (Airflow) / Data Fusion & Tasks / Snowpark orchestration \\
Bulk copy & Copy activity (90+ connectors) & Glue ETL / DMS & Data Fusion / Datastream & \texttt{COPY INTO} / connectors \\
Lookup/enumeration & Lookup, Get Metadata & DynamoDB/S3 lookups in Step Functions & Airflow operators & Metadata queries in tasks \\
Control flow & ForEach, If, Until, Switch & Step Functions states & Airflow DAG logic & Task DAGs \\
Parameterization & Parameters, variables, expressions & Step Functions JSONPath & Airflow templating (Jinja) & Task parameters \\
Code-first alternative & Notebooks in pipelines & Lambda / Glue scripts & Composer DAGs as code & Snowpark Python \\
\bottomrule
\end{tabular}}

\vspace{1mm}
Databricks orchestrates with Workflows; MongoDB with Atlas Triggers and Functions. The portable skill is the metadata-driven mindset this chapter teaches: config rows enumerate work, and one parameterized pipeline executes all of it---that idea survives any tool migration.
\end{crosscloud}

\begin{keyterms}
\textbf{Pipeline}: an orchestrated graph of activities that moves and transforms data on a schedule or trigger. \quad
\textbf{Copy activity}: the managed data-movement activity connecting sources to sinks across dozens of connectors. \quad
\textbf{Lookup}: an activity that reads a small result set (a config table, a file listing) for use by downstream activities. \quad
\textbf{ForEach}: a control-flow activity that iterates child activities over an array. \quad
\textbf{Dynamic expression}: the \texttt{@function(...)} expression language that parameterizes activity properties at run time. \quad
\textbf{Data contract}: a producer-published, consumer-validated specification of schema and semantics for a dataset.
\end{keyterms}

\section{Week 9 Roadmap}
Week~9 opens Part~III, the batch-processing and orchestration arc. Monday introduces Fabric pipelines: Copy, Lookup, and Get Metadata. Tuesday covers control flow---ForEach, If Condition, Until, Switch. Wednesday is the conceptual core: parameters, variables, and dynamic expressions, plus data contracts and schema enforcement. Thursday is the hands-on project: a metadata-driven pipeline that reads table names from a configuration source and copies each from Azure SQL Database to OneLake, with an expected-schema validation gate. Friday scales the pattern into an ingestion framework that onboards hundreds of tables automatically.

Pipelines are the circulatory system of a data platform. Chapter~5's \texttt{COPY INTO} loads files already staged; pipelines are how the files get staged, on a schedule, with retries and alerts \cite{microsoftFabricPipelines}. The metadata-driven pattern taught this week is the difference between an estate of 300 hand-built pipelines and one parameterized pipeline governed by a config table.

\begin{notebox}
Fabric Pipelines inherits the Azure Data Factory lineage: if you have ADF or Synapse pipeline experience, the concepts transfer directly. What changes in Fabric is the destination gravity---everything lands in OneLake---and the shared capacity billing model.
\end{notebox}

\section{Monday: Pipelines, Copy, Lookup, and Get Metadata}
\index{Copy activity}\index{Lookup activity}\index{Get Metadata}
A pipeline is a directed graph of \emph{activities}. Three activities dominate ingestion work \cite{microsoftFabricPipelines,microsoftFabricDataFactory}:

\begin{itemize}
    \item \textbf{Copy activity} moves data between a source and a sink through managed connectors: databases, SaaS applications, file stores, and OneLake itself. It supports full and incremental loads, schema mapping, and fault tolerance settings.
    \item \textbf{Lookup} executes a query or reads a small file and returns its result to the pipeline---classically, ``which tables are active in the config table right now?''
    \item \textbf{Get Metadata} inspects a storage location or dataset---file existence, column lists, modified dates---and returns structural facts the pipeline can branch on.
\end{itemize}

\begin{definitionbox}
A \textbf{metadata-driven pipeline} is one whose work items come from data---a configuration table or manifest---rather than from activities hardcoded per source. Adding a new source is an \texttt{INSERT}, not a deployment.
\end{definitionbox}

The mental model: Lookup and Get Metadata \emph{decide what to do}; Copy \emph{does it}; control flow (Tuesday) \emph{shapes the doing}; parameters and expressions (Wednesday) \emph{make it reusable}.

\begin{pitfall}
Lookup is for small control sets, not data movement. A Lookup that returns a million rows to a ForEach creates a million Copy activities---slow, expensive, and un-debuggable. Enumerate \emph{batches} with Lookup; move \emph{rows} with Copy.
\end{pitfall}

\section{Tuesday: Control Flow}
\index{ForEach}\index{If Condition}\index{Until}\index{Switch}
Control-flow activities shape execution \cite{microsoftFabricPipelines}:

\begin{table}[H]
\centering
\small
\begin{tabular}{@{}p{2.8cm}p{5.6cm}p{5.4cm}@{}}
\toprule
\textbf{Activity} & \textbf{Behavior} & \textbf{Canonical use} \\
\midrule
ForEach & Iterates child activities over an array; sequential by default, parallel with a batch-size cap & Copy each table returned by a Lookup \\
If Condition & Branches on a boolean expression & Route breaking schema changes to quarantine \\
Until & Loops until a condition holds or timeout & Poll an external job's completion status \\
Switch & Multi-way branch on an expression's value & Dispatch per-source-type loading logic \\
\bottomrule
\end{tabular}
\caption{Control-flow activities and their canonical uses.}
\label{tab:chapter9-control-flow}
\end{table}

Every activity also carries dependency conditions---\emph{On Success}, \emph{On Failure}, \emph{On Completion}, \emph{On Skip}---that wire activities into a graph rather than a line. Chapter~12 builds the full orchestration and alerting discipline on these edges; this week's discipline is simpler: every Copy inside a loop should have an explicit failure path, even if that path is only ``record the failure and continue.''

\begin{tipbox}
ForEach runs sequentially by default. For many-table ingestion, enable parallelism with a bounded batch count (start at 4--8). Unbounded parallelism against a source database is a self-inflicted denial-of-service attack.
\end{tipbox}

\section{Wednesday: Parameters, Variables, and Data Contracts}
\index{parameters}\index{variables}\index{dynamic expressions}\index{data contracts}\index{schema enforcement}
\subsection{Parameters and Variables}
\emph{Parameters} are set at trigger time and immutable during the run---use them for batch dates, environment names, and source-system selectors. \emph{Variables} mutate during the run---use them for watermarks advanced mid-flight, accumulated counts, and control flags. \emph{Dynamic expressions} (\texttt{@concat(...)}, \texttt{@activity('Lookup').output.value}, \texttt{@item().table\_name}) wire the two into activity properties at run time.

\subsection{Data Contracts}
A data contract is the producer's published promise about a dataset's schema and semantics, and the consumer's right to validate against it. In Fabric the natural split is: the \emph{producer workspace} publishes and versions the contract; \emph{consumer workspaces} validate incoming batches against it before writing to Silver \cite{microsoftFabricPipelines}.

Enforcement has two layers. At the storage layer, Delta Lake enforces schema by default---writes with mismatched columns are rejected---while \texttt{mergeSchema} and \texttt{overwriteSchema} permit controlled evolution \cite{deltaLakeDocs}. At the pipeline layer, an explicit validation step classifies each incoming change:

\begin{table}[H]
\centering
\small
\begin{tabular}{@{}p{3.4cm}p{4.8cm}p{5.6cm}@{}}
\toprule
\textbf{Change} & \textbf{Class} & \textbf{Pipeline response} \\
\midrule
New nullable column added & Non-breaking (additive) & Allow; log contract version bump \\
Column renamed or dropped & Breaking & Reject batch; quarantine; alert producer \\
Column type widened compatibly & Borderline & Allow only if contract versioning policy says so \\
Required column loses nullability & Breaking & Reject batch; alert producer \\
\bottomrule
\end{tabular}
\caption{Classifying schema changes against a data contract.}
\label{tab:chapter9-contracts}
\end{table}

\begin{pitfall}
Delta's default schema enforcement catches structural violations but not \emph{semantic} ones. A renamed column arriving with the same type as an existing column can pass storage enforcement while corrupting meaning. The contract validation step---an explicit column-list comparison---is what catches it.
\end{pitfall}

\section{Thursday: Hands-on Project}
\index{hands-on lab}\index{metadata-driven ingestion!hands-on}
The Week~9 lab builds the pattern this part of the book is named for: a dynamic pipeline that looks up a list of table names from a configuration source and uses a ForEach loop to copy each from an Azure SQL Database to OneLake---plus an expected-schema validation step that rejects breaking changes and allows additive ones.

\subsection{Goal and Architecture}
Create workspace \texttt{fabric-de-week09} with lakehouse \texttt{lh\_ingest}. Stand up (or simulate with a small Azure SQL free tier) a source database containing \texttt{cfg.ingest\_config} and three sample tables. The pipeline \texttt{PL\_Metadata\_Ingest} reads the config, validates each table's schema against a recorded expectation, and copies conforming tables to Bronze Delta tables partitioned by batch tag.

\subsection{Step-by-Step Actions}
\begin{enumerate}
    \item Create \texttt{cfg.ingest\_config} with columns \texttt{schema\_name}, \texttt{table\_name}, \texttt{is\_active}, \texttt{watermark}, and \texttt{expected\_columns} (a delimited column list).
    \item Insert three rows for the sample tables; seed each sample table with a few hundred rows.
    \item Build the pipeline: Lookup over active config rows feeding a ForEach.
    \item Inside the ForEach, add a script/lookup step that compares the source's current column list to \texttt{expected\_columns}.
    \item Branch with an If Condition: additive-only differences proceed to Copy; renamed or dropped columns route to a Set Variable that records the rejection.
    \item Copy each valid table to \texttt{bronze.<table\_name>} with a batch-tag partition.
    \item Test: rename a source column (batch rejected), add a nullable column (batch allowed), and confirm both outcomes are recorded.
\end{enumerate}

\begin{lstlisting}[language=json,caption={Pipeline JSON (abridged): Lookup feeds ForEach; expressions parameterize the Copy.},label={lst:chapter9-pipeline-json}]
{
  "name": "PL_Metadata_Ingest",
  "variables": { "vBatchTag": { "type": "String", "defaultValue": "daily" } },
  "activities": [
    { "name": "LookupTables", "type": "Lookup",
      "typeProperties": { "sqlReaderQuery":
        "SELECT schema_name, table_name FROM cfg.ingest_config WHERE is_active = 1" } },
    { "name": "CopyEachTable", "type": "ForEach",
      "dependsOn": [ { "activity": "LookupTables", "dependencyConditions": ["Succeeded"] } ],
      "typeProperties": {
        "items": "@activity('LookupTables').output.value",
        "activities": [ { "name": "CopyToOneLake", "type": "Copy",
          "typeProperties": {
            "source": { "sqlReaderQuery":
              "@concat('SELECT * FROM ', item().schema_name, '.', item().table_name)" },
            "sink": { "type": "LakehouseTable",
              "tableOption": "autoCreate",
              "partitionOption": "@concat(variables('vBatchTag'), '/',
                                          item().table_name)" } } } ] } } ]
}
\end{lstlisting}

\begin{lstlisting}[language=Python,caption={Contract validation as a notebook gate: breaking changes reject the batch.},label={lst:chapter9-contract-gate}]
# Notebook invoked per table before the Copy (or as post-landing gate)
expected = ["order_id", "order_date", "customer_id", "net_amount"]
actual = [f.name for f in spark.read.table("bronze.orders_staging").schema]

missing = set(expected) - set(actual)   # renamed/dropped -> breaking
added   = set(actual) - set(expected)   # new columns      -> additive

assert not missing, f"BREAKING contract change, missing columns: {missing}"
if added:
    print(f"NON-BREAKING additive columns accepted: {sorted(added)}")
\end{lstlisting}

\subsection{Validation Checklist}
\begin{itemize}
    \item The pipeline copies all active config rows with no per-table authoring.
    \item A renamed source column rejects the batch and records the rejection.
    \item A new nullable source column passes and is logged as additive.
    \item Adding a fourth table to the config requires zero pipeline edits.
    \item Bronze tables carry the batch-tag partition column.
    \item ForEach parallelism, if enabled, is bounded and documented.
\end{itemize}

\section{Friday: Use Case and Architecture}
\index{metadata-driven ingestion!framework}\index{ingestion framework}
A wholesale platform acquires two companies per year; each arrives with 50--150 operational tables that must flow to OneLake. Hand-building a pipeline per table is a permanent backlog. The deliverable is a metadata-driven ingestion framework: one parameterized pipeline family governed entirely by configuration.

\subsection{The Framework Anatomy}
\begin{itemize}
    \item \textbf{Configuration store}: \texttt{cfg.ingest\_config} rows per table---source connection key, query or full-load flag, watermark column, expected schema, target layer, owner, SLA.
    \item \textbf{Orchestrator}: a master pipeline that looks up due tables (by schedule column) and ForEach-dispatches to a child pipeline, bounded parallelism.
    \item \textbf{Worker}: a parameterized child pipeline executing contract validation, copy, watermark advancement, and run logging for exactly one table.
    \item \textbf{Run ledger}: \texttt{cfg.ingest\_run} rows per attempt---status, rows, duration, error---feeding the Monitoring Hub story of Chapter~12.
    \item \textbf{Onboarding path}: a new table is one config row plus a recorded contract, reviewed like code.
\end{itemize}

\begin{figure}[H]
    \centering
    \begin{tikzpicture}[node distance=8mm and 10mm]
        \node[store] (cfg) at (0,0) {cfg.ingest\_config\\+ run ledger};
        \node[stage] (master) at (4.6,0) {Master pipeline\\(Lookup, ForEach)};
        \node[stage] (worker) at (9.6,0) {Worker pipeline\\(validate, copy, log)};
        \node[store] (lh) at (14.4,0) {OneLake\\Bronze Delta};

        \draw[arrow] (cfg) -- (master);
        \draw[arrow] (master) -- node[above, font=\scriptsize]{per table} (worker);
        \draw[arrow] (worker) -- (lh);
        \draw[arrow, dashed] (worker.south) to[out=-90,in=-90] node[below, font=\scriptsize]{run rows} (cfg.south);

        \node[note, text width=12.5cm, align=center] at (7.2,-2.4) {Onboarding a table = inserting a config row. No pipeline is ever edited to add a source.};
    \end{tikzpicture}
    \caption{Metadata-driven ingestion framework.}
    \label{fig:chapter9-framework}
\end{figure}

\begin{pitfall}
Metadata-driven frameworks amplify both discipline and sloppiness. One bad default in the worker pipeline now affects every table. Treat the worker's contract-validation and run-logging behavior as the most reviewed code in the estate, and version the config schema itself.
\end{pitfall}

\section{Operational Guidance for Week 9}
A pipeline you cannot rerun safely is a liability. Every pattern this week---manifests, watermarks, contract gates---exists so that a failed 3 a.m.\ run can be retried by an on-call engineer following a runbook, not reverse-engineered.

\subsection{Performance and Reliability}
Advance watermarks only on success. Bound ForEach parallelism. Prefer incremental queries (\texttt{WHERE ts > @watermark}) over full reloads once tables outgrow a few million rows. Keep Lookup result sets small. Record every run in the ledger---the ledger is what makes ``did last night work?'' a query rather than an investigation.

\subsection{Security and Governance Checklist}
\begin{itemize}
    \item Store connection credentials in managed connections or Key Vault references, never in pipeline JSON.
    \item Restrict who can edit \texttt{cfg.ingest\_config}---it is executable configuration.
    \item Version the contract per table; record the owning producer workspace.
    \item Quarantine rejected batches to a visible location with reasons, not to the void.
    \item Review worker-pipeline changes with the rigor of a schema migration.
    \item Document each table's SLA and retry unit in the config row.
\end{itemize}

\section{Interview Q\&A}
\begin{enumerate}
    \item \textbf{Q: What makes an ingestion framework metadata-driven?}\\
    A: Work items come from configuration data---a table of sources, schemas, and schedules---that a parameterized pipeline reads at run time, so onboarding a source is an \texttt{INSERT}, not a new pipeline.
    \item \textbf{Q: Which activity iterates over a Lookup result set?}\\
    A: ForEach, with the Lookup's \texttt{output.value} as its \texttt{@items}; parallelism is opt-in and should be bounded.
    \item \textbf{Q: What distinguishes a breaking from a non-breaking schema change?}\\
    A: Breaking changes remove or rename what consumers depend on (dropped/renamed columns, tightened nullability); non-breaking changes only add---a new nullable column. The contract validation gate rejects the first and allows the second.
    \item \textbf{Q: Who owns the data contract: producer or consumer workspace?}\\
    A: The producer publishes and versions it; consumers validate against it. Ownership ambiguity is itself a contract failure.
    \item \textbf{Q: What does \texttt{mergeSchema} allow that default Delta writes reject?}\\
    A: Writes whose schema adds columns beyond the table's current schema; default enforcement rejects any mismatch, which is exactly the protection you want until a reviewed contract bump says otherwise.
    \item \textbf{Q: Why should watermarks advance only on success?}\\
    A: Because the watermark defines what the next run reads. Advancing on failure silently skips uncommitted data; advancing on success makes rerun safe and duplicates impossible.
\end{enumerate}

\section{Further Reading}
Review the Data Factory in Fabric documentation \cite{microsoftFabricDataFactory} and pipeline activity reference \cite{microsoftFabricPipelines}. The schema-enforcement behavior underneath is Delta Lake's \cite{deltaLakeDocs}. For the orchestration layer that monitors these pipelines, Chapter~12 builds on the Monitoring Hub \cite{microsoftFabricMonitoringHub}. Metadata-driven patterns generalize the warehouse loading discipline of Chapter~5 \cite{microsoftFabricCopyInto}.

\section*{Key Takeaways}
\begin{itemize}
    \item Pipelines compose Copy, Lookup, and Get Metadata with control flow to move and validate data on a schedule.
    \item Lookup enumerates work; Copy moves data; never reverse the roles at scale.
    \item Parameters are immutable per run, variables mutate, and dynamic expressions wire both into activities.
    \item Data contracts make schema changes explicit: additive changes flow, breaking changes are rejected and quarantined.
    \item Delta's default schema enforcement is the storage-layer backstop; the pipeline's validation step is the semantic gate.
    \item Metadata-driven design turns onboarding into configuration and concentrates risk in one well-reviewed worker pipeline.
    \item Watermarks advance on success only---that single rule is most of idempotent ingestion.
\end{itemize}

\section{Exercises}
\begin{enumerate}
    \item \textbf{(Conceptual)} Explain why ``enumerate batches, move rows'' is the correct division of labor between Lookup and Copy.
    \item \textbf{(Conceptual)} Classify each as breaking or non-breaking: new nullable column; renamed column; widened \texttt{INT} to \texttt{BIGINT}; column reordered; dropped column.
    \item \textbf{(Hands-on)} Add an Until activity that polls an external job status every 60 seconds with a 20-minute timeout.
    \item \textbf{(Hands-on)} Extend the lab pipeline with a run-ledger write after each Copy---status, rows, duration.
    \item \textbf{(Design)} Design the config schema for a framework spanning Azure SQL, REST APIs, and SFTP sources. Which columns become polymorphic, and how do you keep the worker pipeline sane?
    \item \textbf{(Design)} Two hundred tables, three source systems, one capacity: propose ForEach parallelism settings and scheduling windows, with reasoning.
    \item \textbf{(Interview)} In five sentences, explain how a contract gate and Delta schema enforcement complement rather than duplicate each other.
    \item \textbf{(Operational)} Write the runbook for ``config row added but table never lands''---the five checks in order.
\end{enumerate}

\section{Capstone Project: Config-Driven Multi-Source Ingestion}\label{sec:ch9-capstone}
\index{capstone project}\index{metadata-driven ingestion!capstone}\index{data contracts!capstone}

\begin{capstonebox}
\textbf{Mission.} Build a two-tier, metadata-driven ingestion framework: a master pipeline that dispatches due tables to a parameterized worker pipeline performing contract validation, incremental copy, watermark advancement, and run logging. Prove it by onboarding a new table with a single \texttt{INSERT}, and by rejecting a deliberately breaking schema change.

\textbf{Estimated time.} 5--7 hours.

\textbf{What you'll practice.}
\begin{itemize}
    \item Master/worker pipeline composition with bounded ForEach parallelism.
    \item Contract validation distinguishing breaking from additive changes.
    \item Watermark-driven incremental loads with a success-only advance rule.
    \item Run-ledger design that makes operations queryable.
\end{itemize}
\end{capstonebox}

\subsection*{Scenario}
Contoso's analytics platform must absorb three source systems---an Azure SQL ERP, a REST export dropped as files, and a legacy CSV feed---with new tables arriving monthly. The data engineering team is two people. The framework, not headcount, must absorb the growth.

\subsection*{Requirements}
\begin{itemize}
    \item \texttt{cfg.ingest\_config} and \texttt{cfg.ingest\_run} tables with the columns this chapter specifies.
    \item Master pipeline \texttt{PL\_Master\_Ingest} dispatching due config rows to worker \texttt{PL\_Worker\_Ingest} with parallelism capped at four.
    \item Worker pipeline performing: contract gate $\rightarrow$ incremental copy \texttt{WHERE ts > @watermark} $\rightarrow$ watermark advance on success $\rightarrow$ ledger write on every outcome.
    \item Evidence: one table onboarded by config insert only; one breaking change rejected with quarantine; one additive change flowing through.
    \item A README section describing the onboarding procedure in five steps.
\end{itemize}

\subsection*{Implementation Steps}
\begin{enumerate}
    \item Create the config and ledger tables; seed four tables across two source types.
    \item Build the worker pipeline end-to-end for one table, manually triggered.
    \item Build the master pipeline; verify dispatch and parallelism cap.
    \item Execute the three evidence scenarios and capture ledger rows for each.
    \item Write the onboarding README and a failure runbook.
\end{enumerate}

\subsection*{Deliverables}
\begin{itemize}
    \item Pipeline JSON exports (master and worker) under version control.
    \item Config and ledger table contents after the evidence scenarios.
    \item Screenshot of a bounded-parallelism ForEach run in the Monitoring Hub.
    \item Onboarding README and failure runbook.
\end{itemize}

\subsection*{Acceptance Criteria}
\begin{itemize}
    \item[\(\square\)] A new table flows end-to-end after one config \texttt{INSERT} and no pipeline edits.
    \item[\(\square\)] A renamed source column rejects the batch, writes a quarantine record, and alerts via the ledger.
    \item[\(\square\)] A new nullable source column flows through and is logged as an additive contract change.
    \item[\(\square\)] Watermarks advance only on successful runs; a failed run replays the same slice.
    \item[\(\square\)] Every run---success, rejection, or failure---produces exactly one ledger row.
    \item[\(\square\)] No credentials appear in pipeline JSON or the repository.
\end{itemize}

\subsection*{Stretch Goals}
\begin{itemize}
    \item Add per-table SLA columns and a ledger query that flags tables breaching their SLA.
    \item Add a Switch activity dispatching different copy shapes per source type.
    \item Generate the config row for a new source table automatically from source system metadata.
    \item Publish a small Power BI operations report over the run ledger.
\end{itemize}
